\chapter{پیاده‌سازی}

این فصل پیاده‌سازی را از زاویه رفتار سامانه توضیح می‌دهد. به جای تکیه روی جزئیات داخلی پروژه، مسیر ذهنی کاربر و نقش هر بخش فنی بیان می‌شود؛ چون خواننده گزارش باید بفهمد محصول چگونه کار می‌کند، نه این‌که درگیر نام داخلی فایل‌ها شود.

\section{مسیرهای بک‌اند}
بک‌اند مجموعه‌ای از مسیرهای نسخه‌دار دارد که بخش‌های اصلی محصول را پوشش می‌دهند: سلامت سامانه، قرارداد \lr{OpenAPI}، ثبت‌نام، ورود، خروج، حساب کاربری، پروفایل، \lr{API Key}، سازمان‌ها، اشتراک، مصرف، کاتالوگ، مستندات، \lr{Endpoint}ها، \lr{Caller} و \lr{Studio}. استفاده از مسیرهای نسخه‌دار برای \lr{API} کمک می‌کند توسعه آینده کنترل‌شده‌تر باشد و \lr{Client}ها با تغییرات ناگهانی روبه‌رو نشوند \cite{drf,openapi}.

این مسیرها دو سطح تجربه می‌سازند. سطح اول عمومی است و برای مشاهده کاتالوگ، جزئیات \lr{API} و مستندات استفاده می‌شود. سطح دوم خصوصی است و بعد از ورود کاربر فعال می‌شود؛ مثل \lr{Dashboard}، مصرف، اشتراک، سازمان و انتشار \lr{API}.

\section{لایه \lr{Repository}}
لایه \lr{Repository} منطق کار با داده را متمرکز می‌کند. ساخت کاربر، نگهداری نشست، چرخش \lr{API Key}، خواندن کاتالوگ، ثبت مصرف \lr{Caller} و ذخیره جریان‌های \lr{Studio} از همین مسیر انجام می‌شود. این شیوه باعث می‌شود \lr{View}ها بیش از حد شلوغ نشوند و جزئیات داده در یک لایه مشخص بماند.

در پروژه‌ای شبیه بازارچه \lr{API}، این جداسازی مهم است. داده‌های حساب، کاتالوگ، مستندات، اشتراک و مصرف شکل یکسانی ندارند. وقتی این بخش‌ها پشت یک \lr{Repository} قرار می‌گیرند، تغییر در روش ذخیره‌سازی کمتر به بخش‌های دیگر نشت می‌کند.

\section{\lr{Serializer} و قرارداد پاسخ}
\lr{Serializer}ها نقش مرز امن بین داده داخلی و خروجی قابل نمایش را دارند. داده‌ای که داخل سامانه نگهداری می‌شود، همیشه نباید با همان شکل به کاربر برگردد. برای نمونه، مقدار کامل \lr{API Key} نباید در پاسخ عادی دیده شود و بهتر است فقط نسخه پنهان‌شده آن نمایش داده شود. این رویکرد با نقش \lr{Serializer} در \lr{Django REST Framework} هماهنگ است \cite{drf}.

در بخش انتشار \lr{API} نیز \lr{Serializer} کمک می‌کند داده ورودی قبل از ثبت بررسی شود: نام، آدرس پایه، روش احراز هویت، دسته و توضیح سرویس باید روشن باشند. نتیجه این کار، پاسخ‌هایی است که برای فرانت‌اند قابل پیش‌بینی‌ترند.

\section{فرانت‌اند}
فرانت‌اند تجربه اصلی کاربر را می‌سازد: صفحه مرور \lr{API}ها، جزئیات سرویس، مستندات، قیمت‌گذاری، پرداخت، ورود، ثبت‌نام، \lr{Dashboard}، سازمان‌ها، انتشار، \lr{Caller} و \lr{Studio}. معماری کامپوننتی \lr{React} برای چنین تجربه‌ای مناسب است، چون بخش‌های تکرارشونده مثل فرم، کارت، جدول و وضعیت بارگذاری را می‌توان جدا نگه داشت \cite{react}.

ارتباط فرانت‌اند با بک‌اند از طریق \lr{Client}های مشخص انجام می‌شود. یک گروه درخواست‌ها برای کاتالوگ است، یک گروه برای احراز هویت، و یک گروه برای حساب کاربری و ابزارهای توسعه‌دهنده. این جداسازی باعث می‌شود صفحه‌ها کمتر به جزئیات آدرس‌ها وابسته باشند و نگهداری پروژه ساده‌تر شود.

\section{جریان‌های اصلی پیاده‌سازی}
جدول \ref{tab:implementation-flows} مسیرهای اصلی محصول را از نگاه کاربر و سامانه نشان می‌دهد.

\begin{longtable}{p{3.0cm} p{4.6cm} p{4.8cm}}
\caption{جریان‌های اصلی در پیاده‌سازی}\label{tab:implementation-flows}\\
\toprule
\rowcolor{tablehead}\textbf{جریان} & \textbf{بخش‌های درگیر} & \textbf{خروجی برای کاربر} \\
\midrule
\endfirsthead
\toprule
\rowcolor{tablehead}\textbf{جریان} & \textbf{بخش‌های درگیر} & \textbf{خروجی برای کاربر} \\
\midrule
\endhead
\ModernTableStart
مرور کاتالوگ & فرانت‌اند، \lr{Client} کاتالوگ، بک‌اند و لایه داده & فهرست قابل جستجو و صفحه‌بندی‌شده \lr{API}ها \\
جزئیات \lr{API} & صفحه جزئیات، مستندات، پلن‌ها و \lr{Endpoint}ها & شناخت سرویس قبل از مصرف \\
ورود و نشست & فرم ورود، احراز هویت، کوکی نشست و پاسخ کاربر & دسترسی به بخش خصوصی \\
\lr{Dashboard} & حساب کاربری، مصرف، اشتراک، سازمان و \lr{API Key} & مدیریت وضعیت توسعه‌دهنده \\
انتشار \lr{API} & فرم انتشار، اعتبارسنجی و ثبت در کاتالوگ & اضافه‌شدن سرویس جدید به محصول \\
\lr{Caller} و \lr{Studio} & ابزار تست، جریان کاری و ثبت مصرف & آزمایش و طراحی جریان‌های مرتبط با \lr{API} \\
\ModernTableStop
\bottomrule
\end{longtable}

\section{بخش‌های فنی}
جدول \ref{tab:artifacts} نشان می‌دهد هر بخش فنی چه نقشی در گزارش دارد، بدون این‌که خواننده درگیر نام‌های داخلی پروژه شود.

\begin{longtable}{p{3.3cm} p{4.6cm} p{4.4cm}}
\caption{بخش‌های فنی و نقش آن‌ها در پروژه}\label{tab:artifacts}\\
\toprule
\rowcolor{tablehead}\textbf{بخش} & \textbf{نقش در سامانه} & \textbf{اثر در گزارش} \\
\midrule
\endfirsthead
\toprule
\rowcolor{tablehead}\textbf{بخش} & \textbf{نقش در سامانه} & \textbf{اثر در گزارش} \\
\midrule
\endhead
\ModernTableStart
مسیریابی بک‌اند & تعریف سطح \lr{API} و نسخه‌بندی مسیرها & توضیح قرارداد ارتباطی سامانه \\
\lr{View} و \lr{Serializer} & دریافت درخواست، اعتبارسنجی و ساخت پاسخ & توضیح رفتار حساب، کاتالوگ و ابزارها \\
\lr{Repository} داده & مدیریت داده‌های کاربر، کاتالوگ، مصرف و اشتراک & توضیح مدل عملیاتی داده \\
رابط کاربری & ساخت صفحه‌ها و جریان حرکت کاربر & توضیح تجربه محصول \\
\lr{Client} فرانت‌اند & ارتباط کنترل‌شده با بک‌اند & توضیح اتصال صفحه‌ها به \lr{API} \\
اجرای محلی & بالا آوردن فرانت‌اند و بک‌اند کنار هم & توضیح مسیر ارائه و آزمون دستی \\
آزمون‌ها & کنترل رفتارهای اصلی و خطاهای مهم & پشتوانه بخش ارزیابی \\
\ModernTableStop
\bottomrule
\end{longtable}

\section{نمونه کد}
کد زیر نمونه‌ای کوتاه از منطق پنهان‌کردن داده حساس را نشان می‌دهد.
\begin{latin}
\begin{lstlisting}[language=Python]
SENSITIVE_KEYS = {"password", "api_key", "token", "secret", "authorization"}

def redact_secrets(value):
    if isinstance(value, dict):
        return {
            key: "[REDACTED]" if str(key).lower() in SENSITIVE_KEYS else redact_secrets(item)
            for key, item in value.items()
        }
    return value
\end{lstlisting}
\end{latin}

این قطعه ساده است، اما برای چنین پروژه‌ای مهم است. سامانه با گذرواژه، \lr{API Key}، توکن و اطلاعات احراز هویت سروکار دارد. اگر این داده‌ها بدون کنترل وارد پاسخ‌ها یا لاگ‌ها شوند، امنیت کاربر آسیب می‌بیند. جایگزین‌کردن مقدار حساس با \lr{[REDACTED]} کمک می‌کند خروجی‌ها برای بررسی فنی قابل استفاده باشند، اما رازهای کاربر لو نروند.

\section{جمع‌بندی فصل}
پیاده‌سازی پروژه از چند بخش روشن ساخته شده است: فرانت‌اند برای تجربه کاربر، بک‌اند برای قرارداد \lr{API}، \lr{Serializer} برای کنترل ورودی و خروجی، و \lr{Repository} برای کار با داده. همین ترکیب باعث می‌شود پروژه به یک پلتفرم واقعی اشتراک‌گذاری \lr{API} نزدیک شود.
